\chapter{2015年天津卷（文）}

\section{单选题}

\begin{problem}
已知全集 \(U = \{1, 2, 3, 4, 5, 6\}\)，集合 \(A = \{2, 3, 5\}\)，集合 \(B = \{1, 3, 4, 6\}\)，则集合 \(A \cap \complement_U B =\)（\quad）

\choices
    {\( \{3\} \)}
    {\(\{2, 5\}\)}
    {\(\{1, 4, 6\}\)}
    {\(\{2, 3, 5\}\)}
\end{problem}

\begin{answer}
B
\end{answer}

\begin{solution}
全集中的补集为
\[
\complement_U B=\{2,5\}
\]
因此
\[
A\cap\complement_U B=\{2,3,5\}\cap\{2,5\}=\{2,5\}
\]
所以选 B.
\end{solution}

\begin{problem}
设变量 \(x\)，\(y\) 满足约束条件 \(\begin{cases}
x - 2 \leqslant 0,\\
x - 2y \leqslant 0,\\
x + 2y - 8 \leqslant 0,
\end{cases}\) 则目标函数 \(z = 3x + y\) 的最大值为（\quad）

\choices
    {7}
    {8}
    {9}
    {14}
\end{problem}

\begin{answer}
C
\end{answer}

\begin{solution}
将目标函数配成约束条件的线性组合：
\[
z=3x+y=\frac{5}{2}(x-2)+\frac{1}{2}(x+2y-8)+9\leqslant9
\]
当 \(x=2\)，\(y=3\) 时，三个约束条件均成立，且 \(z=3\times2+3=9\). 因此最大值为 \(9\)，选 C.
\end{solution}

\begin{problem}
阅读程序框图，运行相应的程序，则输出 \(i\) 的值为（\quad）

\bitmapfigure[max width=6.0cm]{img/2015/tianjin_liberal/q3_fig1.png}

\choices
    {2}
    {3}
    {4}
    {5}
\end{problem}

\begin{answer}
C
\end{answer}

\begin{solution}
程序开始时 \(S=10\)，\(i=0\). 每循环一次先令 \(i=i+1\)，再令 \(S=S-i\)，各次循环后的结果为
\((i,S)=(1,9)\)，\((2,7)\)，\((3,4)\)，\((4,0)\).
前三次循环后 \(S>1\)，第四次循环后 \(S=0\leqslant1\)，程序停止，此时输出的 \(i\) 为 \(4\). 所以选 C.

这里按题干所问的变量 \(i\) 确定输出；原流程图终止框写为“输出 \(S\)”，若按该框文字执行则输出 \(S=0\)，与题干及选项不一致. 故本题按题干应选 C.
\end{solution}

\begin{problem}
设 \(x \in \R\)，则 “\(1 < x < 2\)” 是 “\(|x - 2| < 1\)” 的（\quad）

\choices
    {充分而不必要条件}
    {必要而不充分条件}
    {充要条件}
    {既不充分也不必要条件}
\end{problem}

\begin{answer}
A
\end{answer}

\begin{solution}
由 \(1<x<2\) 可得 \(-1<x-2<0\)，从而 \(\lvert x-2\rvert<1\)，所以前一个条件是后一个条件的充分条件.

但取 \(x=\frac{5}{2}\) 时，\(\lvert x-2\rvert=\frac12<1\)，而 \(1<x<2\) 不成立，故不是必要条件. 因此选 A.
\end{solution}

\begin{problem}
已知双曲线 \(\frac{x^2}{a^2} - \frac{y^2}{b^2} = 1\)（\(a > 0\)，\(b > 0\)） 的一焦点为 \(F(2,0)\)，且双曲线的渐近线与圆 \((x - 2)^2 + y^2 = 3\) 相切，则双曲线的方程为（\quad）

\choices
    {\(\frac{x^2}{4} -\frac{y^2}{9} = 1\)}
    {\(\frac{x^2}{13} -\frac{y^2}{9} = 1\)}
    {\(\frac{x^2}{3} - y^2 = 1\)}
    {\(x^{2} - \frac{y^{2}}{3} = 1\)}
\end{problem}

\begin{answer}
D
\end{answer}

\begin{solution}
设双曲线的半焦距为 \(c\)，则由焦点 \(F(2,0)\) 得 \(c=2\)，并且
\[
a^2+b^2=c^2=4
\]
双曲线的渐近线为 \(y=\pm\frac{b}{a}x\). 取其中一条写成 \(bx-ay=0\)，圆心 \((2,0)\) 到该直线的距离为
\[
\frac{2b}{\sqrt{a^2+b^2}}=\frac{2b}{2}=b
\]
由于渐近线与半径为 \(\sqrt3\) 的圆相切，故 \(b=\sqrt3\)，即 \(b^2=3\). 于是 \(a^2=4-3=1\)，双曲线方程为
\[
\frac{x^2}{1}-\frac{y^2}{3}=1
\]
所以选 D.
\end{solution}

\begin{problem}
如图，在圆 \(O\) 中，\(M\)，\(N\) 是弦 \(AB\) 的三等分点，弦 \(CD\)，\(CE\) 分别经过点 \(M\)，\(N\)，若 \(CM = 2\)，\(MD = 4\)，\(CN = 3\)，则线段 \(NE\) 的长为（\quad）

\bitmapfigure[max width=0.36\linewidth]{img/2015/tianjin_liberal/q6_fig1.png}

\choices
    {\(\frac{8}{3}\)}
    {3}
    {\( \frac{10}{3} \)}
    {\( \frac{5}{2} \)}
\end{problem}

\begin{answer}
A
\end{answer}

\begin{solution}
设 \(AM=MN=NB=t\). 由相交弦定理，在点 \(M\) 处有
\[
MA\cdot MB=MC\cdot MD=2\times4=8
\]
而 \(MA=t\)，\(MB=2t\)，所以 \(2t^2=8\)，从而 \(t=2\). 于是 \(NA=NM+MA=4\)，\(NB=2\).

再由相交弦定理，在点 \(N\) 处有
\[
NC\cdot NE=NA\cdot NB=4\times2=8
\]
已知 \(NC=3\)，故 \(NE=\frac{8}{3}\)，选 A.
\end{solution}

\begin{problem}
已知定义在 \(\R\) 上的函数 \(f(x) = 2^{|x - m|} - 1\) （\(m\) 为实数） 为偶函数，记 \(a = f(\log_{0.5}3)\)，\(b = f(\log_2 5)\)，\(c = f(2m)\)，则 \(a\)，\(b\)，\(c\) 的大小关系为（\quad）

\choices
    {\(a < b < c\)}
    {\(c < a < b\)}
    {\(a < c < b\)}
    {\(c < b < a\)}
\end{problem}

\begin{answer}
B
\end{answer}

\begin{solution}
因为 \(f(x)=2^{\lvert x-m\rvert}-1\) 为偶函数，所以对任意 \(x\) 有
\[
2^{\lvert x-m\rvert}=2^{\lvert -x-m\rvert}
\]
指数函数 \(2^t\) 严格递增，故 \(\lvert x-m\rvert=\lvert -x-m\rvert\). 两边平方后得到
\[
(x-m)^2=(-x-m)^2
\]
即 \(4xm=0\). 这对任意 \(x\) 成立，所以 \(m=0\).

于是
\[
a=2^{\lvert\log_{0.5}3\rvert}-1=2^{\log_2 3}-1=2
\]
\[
b=2^{\lvert\log_2 5\rvert}-1=2^{\log_2 5}-1=4
\]
\[
c=2^{\lvert2m-m\rvert}-1=0
\]
因此 \(c<a<b\)，选 B.
\end{solution}

\begin{problem}
已知函数 \( f(x) = \begin{cases}
2 - |x|, & x \leqslant 2,\\
(x - 2)^2, & x > 2,
\end{cases} \) 函数 \( g(x) = 3 - f(2 - x) \)，则函数 \( y = f(x) - g(x) \) 的零点个数为（\quad）

\choices
    {2}
    {3}
    {4}
    {5}
\end{problem}

\begin{answer}
A
\end{answer}

\begin{solution}
由 \(y=f(x)-g(x)=0\) 及 \(g(x)=3-f(2-x)\)，零点满足
\[
f(x)+f(2-x)=3
\]

当 \(x\leqslant0\) 时，\(f(x)=2+x\)，且 \(2-x\geqslant2\)，所以 \(f(2-x)=x^2\). 方程化为
\[
x^2+x-1=0
\]
其根为 \(\frac{-1\pm\sqrt5}{2}\)，其中只有 \(\frac{-1-\sqrt5}{2}\leqslant0\) 符合，得到一个零点.

当 \(0\leqslant x\leqslant2\) 时，\(f(x)=2-x\)，\(f(2-x)=x\)，故 \(f(x)+f(2-x)=2\)，没有零点.

当 \(x>2\) 时，\(f(x)=(x-2)^2\)，而 \(2-x<0\)，故 \(f(2-x)=4-x\). 方程化为
\[
(x-2)^2+4-x=3
\]
即 \(x^2-5x+5=0\). 其根为 \(\frac{5\pm\sqrt5}{2}\)，其中只有 \(\frac{5+\sqrt5}{2}>2\) 符合，得到一个零点.

综上，零点共有 \(2\) 个，选 A.
\end{solution}

\section{填空题}

\begin{problem}
\(\i\) 是虚数单位，计算 \( \frac{1-2\i}{2+\i} \) 的结果为\fillinblank{}.
\end{problem}

\begin{answer}
\(-\i\)
\end{answer}

\begin{solution}
分子分母同乘分母的共轭复数 \(2-\i\)，得
\[
\frac{1-2\i}{2+\i}=\frac{(1-2\i)(2-\i)}{(2+\i)(2-\i)}=\frac{2-5\i+2\i^2}{5}=\frac{-5\i}{5}=-\i
\]
\end{solution}

\begin{problem}
一个几何体的三视图如图所示（单位：\(\mathrm{m}\)），则该几何体的体积为\fillinblank{}\(\mathrm{m}^{3}\).

\bitmapfigure[max width=0.9\linewidth]{img/2015/tianjin_liberal/q10_fig1.png}
\end{problem}

\begin{answer}
\(\frac{8}{3}\pi\)
\end{answer}

\begin{solution}
由三视图可知，该几何体由一个半径为 \(1\)、高为 \(2\) 的圆柱和圆柱两端各一个底面半径为 \(1\)、高为 \(1\) 的圆锥组成. 因此体积为
\[
V=\pi\times1^2\times2+2\times\frac13\pi\times1^2\times1=\frac83\pi
\]
\end{solution}

\begin{problem}
已知函数 \(f(x) = ax \ln x\)，\(x \in (0, +\infty)\)，其中 \( a \) 为实数，\( f'(x) \) 为 \( f(x) \) 的导函数. 若 \( f'(1) = 3 \)，则 \( a \) 的值为\fillinblank{}.
\end{problem}

\begin{answer}
3
\end{answer}

\begin{solution}
\[
f'(x)=a(\ln x+1)
\]
代入 \(x=1\) 并利用 \(\ln1=0\)，得 \(f'(1)=a=3\).
\end{solution}

\begin{problem}
已知 \(a > 0\)，\(b > 0\)，\(ab = 8\)，则当 \(a\) 的值为\fillinblank{} 时，\(\log_2 a \cdot \log_2 (2b)\) 取得最大值.
\end{problem}

\begin{answer}
4
\end{answer}

\begin{solution}
由 \(ab=8\) 且 \(a\)，\(b>0\)，得 \(b=\frac8a\). 令 \(t=\log_2a\)，则
\[
\log_2(2b)=1+\log_2b=1+\log_2 8-\log_2a=4-t
\]
所以目标式为
\[
t(4-t)=4-(t-2)^2\leqslant4
\]
当且仅当 \(t=2\) 时取最大值. 此时 \(a=2^2=4\).
\end{solution}

\begin{problem}
在等腰梯形 \(ABCD\) 中，已知 \(AB \parallel DC\)，\(AB = 2\)，\(BC = 1\)，\(\angle ABC = 60^\circ\). 点 \(E\) 和 \(F\) 分别在线段 \(BC\) 和 \(DC\) 上，且 \(\overrightarrow{BE} = \frac{2}{3} \overrightarrow{BC}\)，\(\overrightarrow{DF} = \frac{1}{6} \overrightarrow{DC}\)，则 \(\overrightarrow{AE} \cdot \overrightarrow{AF}\) 的值为\fillinblank{}.
\end{problem}

\begin{answer}
\(\frac{29}{18}\)
\end{answer}

\begin{solution}
建立平面直角坐标系，取 \(B=(0,0)\)，\(A=(2,0)\). 由 \(BC=1\) 且 \(\angle ABC=60^\circ\)，可取
\[
C=\left(\frac12,\frac{\sqrt3}{2}\right)
\]
等腰梯形的另一腰等于 \(BC\)，结合 \(AB\parallel DC\) 得
\[
D=\left(\frac32,\frac{\sqrt3}{2}\right)
\]
由分点条件
\[
E=B+\frac23(C-B)=\left(\frac13,\frac{\sqrt3}{3}\right)
\]
\[
F=D+\frac16(C-D)=\left(\frac43,\frac{\sqrt3}{2}\right)
\]
因此
\(\overrightarrow{AE}=\left(-\frac53,\frac{\sqrt3}{3}\right)\)，\(\overrightarrow{AF}=\left(-\frac23,\frac{\sqrt3}{2}\right)\)
从而
\[
\overrightarrow{AE}\cdot\overrightarrow{AF}=\frac{10}{9}+\frac{3}{6}=\frac{29}{18}
\]
\end{solution}

\begin{problem}
已知函数 \(f(x) = \sin \omega x + \cos \omega x\)（\(\omega > 0\)），\( x \in \R \). 若函数 \( f(x) \) 在区间 \((- \omega, \omega)\) 内单调递增，且函数 \( y = f(x) \) 的图象关于直线 \( x = \omega \) 对称，则 \( \omega \) 的值为\fillinblank{}.
\end{problem}

\begin{answer}
\(\frac{\sqrt{\pi}}{2}\)
\end{answer}

\begin{solution}
将函数写成
\[
f(x)=\sqrt{2}\sin\left(\omega x+\frac{\pi}{4}\right)
\]
图象关于直线 \(x=\omega\) 对称，故对任意 \(t\) 有 \(f(\omega+t)=f(\omega-t)\)，于是
\[
2\sqrt{2}\cos\left(\omega^2+\frac{\pi}{4}\right)\sin(\omega t)=0
\]
由于 \(\omega>0\)，这对任意 \(t\) 成立只能是
\[
\cos\left(\omega^2+\frac{\pi}{4}\right)=0
\]
即
\(\omega^2=\frac{(4k+1)\pi}{4}\)，\(k\in\Z\).

又
\[
f'(x)=\omega\bigl(\cos\omega x-\sin\omega x\bigr)=\sqrt{2}\omega\cos\left(\omega x+\frac{\pi}{4}\right)
\]
函数在 \((-\omega,\omega)\) 内单调递增，要求上式在该区间内为正. 若 \(k\leqslant-1\)，则 \(\omega^2<0\) 不可能. 若 \(k\geqslant1\)，则 \(\omega^2\geqslant\frac{5\pi}{4}\)，取 \(x=\frac{3\pi}{4\omega}\)，有 \(|x|<\omega\) 且
\[
f'(x)=\sqrt{2}\omega\cos\pi<0
\]
也不满足单调递增. 因此只能 \(k=0\). 此时 \(\omega^2=\frac{\pi}{4}\)，且当 \(-\omega<x<\omega\) 时
\[
0<\omega x+\frac{\pi}{4}<\frac{\pi}{2}
\]
从而 \(f'(x)>0\)，条件确实满足. 故
\[
\omega=\frac{\sqrt{\pi}}{2}
\]
\end{solution}

\section{解答题}

\begin{problem}
设甲，乙，丙三个乒乓球协会的运动员人数分别为 27， 9， 18. 现采用分层抽样的方法从这三个协会中抽取 6 名运动员组队参加比赛.
\begin{enumerate}
\item 求应从这三个协会中分别抽取的运动员的人数；
\item 将抽取的 6 名运动员进行编号，编号分别为 \(A_{1}\)，\(A_{2}\)，\(A_{3}\)，\(A_{4}\)，\(A_{5}\)，\(A_{6}\). 现从这 6 名运动员中随机抽取 2 人参加双打比赛.
\begin{enumerate}
\item 用所给编号列出所有可能的结果；
\item 设 \(A\) 为事件：“编号为 \(A_{5}\) 和 \(A_{6}\) 的两名运动员中至少有 1 人被抽到”，求事件 \(A\) 发生的概率.
\end{enumerate}
\end{enumerate}
\end{problem}

\begin{solution}
\begin{enumerate}
\item 三个协会共有 \(27+9+18=54\) 名运动员，抽取比例为 \(6/54=1/9\)，所以应分别抽取
\(27\times\frac19=3\)，\(9\times\frac19=1\)，\(18\times\frac19=2\)
人.

\item
\begin{enumerate}
\item 随机抽取的两人不计顺序，所有可能结果为
\[
\begin{gathered}
(A_1,A_2),(A_1,A_3),(A_1,A_4),(A_1,A_5),(A_1,A_6),\\
(A_2,A_3),(A_2,A_4),(A_2,A_5),(A_2,A_6),\\
(A_3,A_4),(A_3,A_5),(A_3,A_6),\\
(A_4,A_5),(A_4,A_6),(A_5,A_6).
\end{gathered}
\]
共 \(\mathrm C_6^2=15\) 种.

\item 事件 \(A\) 的对立事件是 \(A_5\)，\(A_6\) 均未被抽到. 其结果数为从 \(A_1\)，\(A_2\)，\(A_3\)，\(A_4\) 中抽取两人的结果数，即 \(\mathrm C_4^2=6\). 因此事件 \(A\) 的结果数为 \(15-6=9\)，从而
\[
P(A)=\frac9{15}=\frac35
\]
\end{enumerate}
\end{enumerate}
\end{solution}

\begin{problem}
在 \(\triangle ABC\) 中，内角 \(A\)，\(B\)，\(C\) 所对的边分别为 \(a\)，\(b\)，\(c\). 已知 \(\triangle ABC\) 的面积为 \(3\sqrt{15}\)，\(b - c = 2\)，\(\cos A = -\frac{1}{4}\).
\begin{enumerate}
\item 求 \(a\) 和 \( \sin C \) 的值；
\item 求 \( \cos\left(2A+\frac{\pi}{6}\right) \) 的值.
\end{enumerate}
\end{problem}

\begin{solution}
\begin{enumerate}
\item 因为 \(\cos A=-\frac14\)，且 \(0<A<\pi\)，所以
\[
\sin A=\sqrt{1-\cos^2A}=\frac{\sqrt{15}}4
\]
由三角形面积公式
\[
\frac12bc\sin A=3\sqrt{15}
\]
得 \(bc=24\). 结合 \(b-c=2\)，有 \(b=c+2\)，所以
\[
c(c+2)=24
\]
由边长为正，得 \(c=4\)，从而 \(b=6\). 由余弦定理
\[
a^2=b^2+c^2-2bc\cos A=36+16+12=64
\]
故 \(a=8\). 再由正弦定理
\[
\frac{\sin C}{c}=\frac{\sin A}{a}
\]
得
\[
\sin C=\frac{4\cdot\frac{\sqrt{15}}4}{8}=\frac{\sqrt{15}}8
\]

\item
\(\cos2A=2\cos^2A-1=-\frac78\)，\(\sin2A=2\sin A\cos A=-\frac{\sqrt{15}}8\).
因此
\[
\cos\left(2A+\frac\pi6\right)=\cos2A\cos\frac\pi6-\sin2A\sin\frac\pi6
=-\frac{7\sqrt3}{16}+\frac{\sqrt{15}}{16}
=\frac{\sqrt{15}-7\sqrt3}{16}
\]
\end{enumerate}
\end{solution}

\begin{problem}
如图，已知 \(AA_{1} \perp\) 平面 \(ABC\)，\(BB_{1} \parallel AA_{1}\)，\(AB = AC = 3\)，\(BC = 2\sqrt{5}\)，\(AA_{1} = \sqrt{7}\)，\(BB_{1} = 2\sqrt{7}\)，点 \(E\) 和 \(F\) 分别为 \(BC\) 和 \(A_{1}C\) 的中点.

\begin{enumerate}
\item 求证： \(EF \parallel\) 平面 \(A_{1}B_{1}BA\)；
\item 求证：平面 \( AEA_{1} \perp \) 平面 \(BCB_{1}\)；
\item 求直线 \(A_{1}B_{1}\) 与平面 \(BCB_{1}\) 所成角的大小.
\end{enumerate}

\bitmapfigure[max width=0.56\linewidth]{img/2015/tianjin_liberal/q17_fig1.png}
\end{problem}

\begin{solution}
\begin{enumerate}
\item 在三角形 \(BCA_1\) 中，\(E\)，\(F\) 分别是 \(BC\)，\(A_1C\) 的中点，所以由中位线定理
\[
EF\parallel BA_1
\]
又 \(BA_1\subset\) 平面 \(A_1B_1BA\). 平面 \(A_1B_1BA\) 与平面 \(ABC\) 的交线为 \(AB\)，而 \(E\) 是 \(BC\) 的中点且不在 \(AB\) 上，所以 \(E\) 不在平面 \(A_1B_1BA\) 内. 因此
\[
EF\parallel\text{平面 }A_1B_1BA
\]

\item 因为 \(AB=AC\)，且 \(E\) 是 \(BC\) 的中点，所以 \(AE\perp BC\). 过 \(E\) 作直线 \(l\parallel BB_1\). 由于 \(E\in BC\)，且 \(BC\) 与 \(BB_1\) 都在平面 \(BCB_1\) 内，所以 \(l\) 也在平面 \(BCB_1\) 内. 又 \(BB_1\parallel AA_1\)，而 \(AA_1\perp\) 平面 \(ABC\)，故 \(AE\perp l\). 直线 \(BC\)，\(l\) 在点 \(E\) 相交，且都在平面 \(BCB_1\) 内，因此
\[
AE\perp\text{平面 }BCB_1
\]
由于 \(AE\subset\) 平面 \(AEA_1\)，所以
\[
\text{平面 }AEA_1\perp\text{平面 }BCB_1
\]

\item 取 \(N\) 为 \(BC\) 的中点. 在 \(N\) 处作 \(NN_1\parallel AA_1\)，且 \(NN_1=AA_1\). 则 \(N_1\) 在平面 \(BCB_1\) 内，且 \(A_1N_1\parallel AN\). 由等腰三角形性质，\(AN\perp BC\). 过 \(N_1\) 作平面 \(BCB_1\) 内分别平行于 \(BC\)，\(BB_1\) 的两条直线；由于 \(A_1N_1\parallel AN\)，且 \(BB_1\parallel AA_1\perp\) 平面 \(ABC\)，有 \(A_1N_1\perp BC\) 和 \(A_1N_1\perp BB_1\). 因此 \(A_1N_1\perp\) 平面 \(BCB_1\). 于是 \(N_1B_1\) 是直线 \(A_1B_1\) 在平面 \(BCB_1\) 内的射影.

由 \(AB=3\)，\(AA_1=\sqrt7\)，\(BB_1=2\sqrt7\)，得
\[
A_1B_1=\sqrt{AB^2+(BB_1-AA_1)^2}=\sqrt{9+7}=4
\]
又 \(NB=\frac12BC=\sqrt5\)，\(B_1N_1=BB_1-NN_1=\sqrt7\)，所以
\[
N_1B_1=\sqrt{NB^2+B_1N_1^2}=\sqrt{5+7}=2\sqrt3
\]
设直线 \(A_1B_1\) 与平面 \(BCB_1\) 所成角为 \(\theta\)，则
\[
\cos\theta=\frac{N_1B_1}{A_1B_1}=\frac{2\sqrt3}{4}=\frac{\sqrt3}{2}
\]
由于 \(0\leqslant\theta\leqslant\frac\pi2\)，故 \(\theta=30^\circ\).
\end{enumerate}
\end{solution}

\begin{problem}
已知数列 \(\{a_{n}\}\) 是各项均为正数的等比数列，\(\{b_{n}\}\) 是等差数列，且 \(a_{1} = b_{1} = 1\)，\(b_{2} + b_{3} = 2a_{3}\)，\(a_{5} - 3b_{2} = 7\).
\begin{enumerate}
\item 求 \( \{a_{n}\} \) 和 \( \{b_{n}\} \) 的通项公式；
\item 设 \(c_{n} = a_{n}b_{n}\)，\(n \in \N^{*}\)，求数列 \(\{c_{n}\}\) 的前 \(n\) 项和.
\end{enumerate}
\end{problem}

\begin{solution}
\begin{enumerate}
\item 设等比数列 \(\{a_n\}\) 的公比为 \(q>0\)，等差数列 \(\{b_n\}\) 的公差为 \(d\). 由 \(a_1=b_1=1\)，有
\(a_n=q^{n-1}\)，\(b_n=1+(n-1)d\).
条件 \(b_2+b_3=2a_3\) 给出
\[
(1+d)+(1+2d)=2q^2
\]
即 \(3d=2q^2-2\). 条件 \(a_5-3b_2=7\) 给出
\[
q^4-3(1+d)=7
\]
代入 \(3d=2q^2-2\)，得
\(q^4-2q^2-8=0\)，\((q^2-4)(q^2+2)=0\).
由于 \(q>0\)，所以 \(q=2\)，进而 \(d=2\). 因此
\(a_n=2^{n-1}\)，\(b_n=2n-1\).

\item 又 \(c_n=a_nb_n=(2n-1)2^{n-1}\). 令
\[
T_k=(2k-3)2^k
\]
则
\[
T_k-T_{k-1}=(2k-1)2^{k-1}=c_k
\]
所以
\[
\sum_{k=1}^n c_k=T_n-T_0=(2n-3)2^n-(-3)=(2n-3)2^n+3
\]
\end{enumerate}
\end{solution}

\begin{problem}
已知椭圆 \(\frac{x^2}{a^2} + \frac{y^2}{b^2} = 1\)（\(a > b > 0\)） 的上顶点为 \(B\)，左焦点为 \(F\)，离心率为 \(\frac{\sqrt{5}}{5}\).
\begin{enumerate}
\item 求直线 \(BF\) 的斜率；
\item 设直线 \(BF\) 与椭圆交于点 \(P\)（\(P\) 异于点 \(B\)），过点 \(B\) 且垂直于 \(BP\) 的直线与椭圆交于点 \(Q\)（\(Q\) 异于点 \(B\)），直线 \(PQ\) 与 \(y\) 轴交于点 \(M\)，\( |PM| = \lambda |MQ| \).
\begin{enumerate}
\item 求 \( \lambda \) 的值；
\item 若 \( |PM| \sin \angle BQP = \frac{7\sqrt{5}}{9} \)，求椭圆的方程.
\end{enumerate}
\end{enumerate}
\end{problem}

\begin{solution}
\begin{enumerate}
\item 设椭圆半焦距为 \(c\). 由离心率条件
\[
\frac ca=\frac{\sqrt5}{5}=\frac1{\sqrt5}
\]
得 \(a=\sqrt5c\)，再由 \(b^2=a^2-c^2\) 得 \(b=2c\). 于是 \(B=(0,2c)\)，\(F=(-c,0)\)，所以
\[
k_{BF}=\frac{2c}{c}=2
\]

\item
\begin{enumerate}
\item 直线 \(BF\) 的方程为 \(y=2x+2c\). 代入椭圆方程
\[
\frac{x^2}{5c^2}+\frac{y^2}{4c^2}=1
\]
得
\[
\frac{x^2}{5c^2}+\frac{(x+c)^2}{c^2}=1
\]
除去对应于 \(B\) 的交点后，得到
\[
P=\left(-\frac{5c}{3},-\frac{4c}{3}\right)
\]
由于 \(BP\) 的斜率为 \(2\)，过 \(B\) 且垂直于 \(BP\) 的直线为
\[
y=2c-\frac{x}{2}
\]
联立椭圆方程，除去 \(B\) 后得到
\[
Q=\left(\frac{40c}{21},\frac{22c}{21}\right)
\]
直线 \(PQ\) 的斜率为
\[
\frac{\frac{22c}{21}+\frac{4c}{3}}{\frac{40c}{21}+\frac{5c}{3}}=\frac23
\]
故其与 \(y\) 轴交点为 \(M=(0,-\frac{2c}{9})\). 由于 \(P\)，\(M\)，\(Q\) 共线，按横坐标比例
\[
\lambda=\frac{|PM|}{|MQ|}=\frac{\frac{5c}{3}}{\frac{40c}{21}}=\frac78
\]

\item 由上面的坐标
\[
|PM|=\sqrt{\left(\frac{5c}{3}\right)^2+\left(\frac{10c}{9}\right)^2}=\frac{5\sqrt{13}}9c
\]
在点 \(Q\) 处，向量 \(\overrightarrow{QB}\) 与 \(\overrightarrow{QP}\) 的方向向量可分别取 \((-2,1)\) 与 \((-3,-2)\)，故
\[
\sin\angle BQP=\frac{|(-2)(-2)-1(-3)|}{\sqrt5\sqrt{13}}=\frac7{\sqrt{65}}
\]
代入题设条件，得
\[
\frac{5\sqrt{13}}9c\cdot\frac7{\sqrt{65}}=\frac{7\sqrt5}{9}
\]
从而 \(c=1\). 因此 \(a=\sqrt5\)，\(b=2\)，椭圆方程为
\[
\frac{x^2}{5}+\frac{y^2}{4}=1
\]
\end{enumerate}
\end{enumerate}
\end{solution}

\begin{problem}
已知函数 \(f(x) = 4x - x^4\)，\(x \in \R\).
\begin{enumerate}
\item 求 \( f(x) \) 的单调区间；
\item 设曲线 \( y = f(x) \) 与 \( x \) 轴正半轴的交点为 \(P\)，曲线在点 \(P\) 处的切线方程为 \( y = g(x) \). 求证：对于任意的实数 \( x \)，都有 \( f(x) \leqslant g(x) \)；
\item 若方程 \( f(x) = a \) （\( a \) 为实数） 有两个正实数根 \(x_1\)，\(x_2\)，且 \( x_1 < x_2 \)，求证：\( x_2 - x_1 \leqslant -\frac{a}{3} + 4^{\frac{1}{3}} \).
\end{enumerate}
\end{problem}

\begin{solution}
\begin{enumerate}
\item
\[
f'(x)=4-4x^3=4(1-x^3)
\]
当 \(x<1\) 时 \(f'(x)>0\)，当 \(x>1\) 时 \(f'(x)<0\). 所以 \(f(x)\) 在 \((-\infty,1)\) 上单调递增，在 \((1,+\infty)\) 上单调递减.

\item 令 \(r=4^{1/3}\). 曲线与 \(x\) 轴正半轴的交点满足 \(4x-x^4=0\)，所以 \(P=(r,0)\). 又
\[
f'(r)=4-4r^3=-12
\]
故切线为
\[
g(x)=-12(x-r)=-12x+12r
\]
于是
\[
g(x)-f(x)=x^4-16x+12r
\]
利用 \(r^3=4\) 因式分解得
\[
x^4-16x+12r=(x-r)^2(x^2+2rx+3r^2)
=(x-r)^2\bigl((x+r)^2+2r^2\bigr)\geqslant0
\]
所以对于任意的实数 \(x\)，都有 \(f(x)\leqslant g(x)\).

\item 由第（1）问，\(f\) 在 \((0,1)\) 上递增，在 \((1,r)\) 上递减，且 \(f(0)=f(r)=0\)，\(f(1)=3\). 题设有两个不同的正根，故 \(0<a<3\)，并且
\[
0<x_1<1<x_2<r
\]
令 \(u=x_1\)，\(v=x_2\). 由 \(f(u)=f(v)\) 且 \(u<v\)，得
\(4(u-v)=u^4-v^4\)，\(u^3+u^2v+uv^2+v^3=4\).
记 \(r^3=4\)，于是
\[
r^3-v^3=u(u^2+uv+v^2)
\]
因此
\[
r-v=\frac{u(u^2+uv+v^2)}{r^2+rv+v^2}
\]
又由
\[
(u+v)^3\geqslant(u+v)(u^2+v^2)=4=r^3
\]
知 \(r\leqslant u+v\)，从而
\[
r^2+rv+v^2\leqslant(u+v)^2+(u+v)v+v^2
\leqslant3(u^2+uv+v^2)
\]
故 \(r-v\geqslant u/3\)，于是
\[
v-u\leqslant r-\frac{4u}{3}
\]
另一方面 \(a=f(u)=4u-u^4\leqslant4u\)，所以
\[
r-\frac{4u}{3}\leqslant r-\frac a3
\]
代回 \(u=x_1\)，\(v=x_2\)，\(r=4^{1/3}\)，得到
\[
x_2-x_1\leqslant-\frac a3+4^{1/3}
\]
\end{enumerate}
\end{solution}
